% from ega3-1-fr.tex
\oldpage[0\textsubscript{III}]{65}
un système projectif d'objets de $\mathcal C$ dont l'ensemble d'indices $I$
est filtrant à droite~; soient $A=\varprojlim A_\alpha$ et, pour tout
$\alpha\in I$, soit $f_\alpha:A\to A_\alpha$ le morphisme canonique. Pour
tout $\alpha\in I$, les $f_{\alpha\beta}(A_\beta)$ pour
$\alpha\leq\beta$ forment une famille filtrante décroissante de sous-objets
de $A_\alpha$~; le sous-objet
\[
A'_\alpha=\inf_{\beta\geq\alpha}f_{\alpha\beta}(A_\beta)
\]
est dit sous-objet des «~images universelles~» dans $A_\alpha$~; il est
clair que $f_\alpha(A)\subset A'_\alpha$ et
$f_{\alpha\beta}(A'_\beta)\subset A'_\alpha$ pour
$\alpha\leq\beta$~; donc
$(A'_\alpha,f_{\alpha\beta}|A'_\beta)$ est un système projectif et
$A=\varprojlim A'_\alpha$.
\end{env}

\begin{env}[13.1.2]
\label{0.13.1.2-fr}
Étant donné un système projectif $(A_\alpha,f_{\alpha\beta})$ dans
$\mathcal C$, on appelle \emph{condition de Mittag-Leffler} la condition
suivante~:
\begin{itemize}
\item[\textbf{(ML)}]
\emph{Pour tout indice $\alpha$, il existe $\beta\geq\alpha$ tel que, pour
tout $\gamma\geq\beta$, on ait
$f_{\alpha\gamma}(A_\gamma)=f_{\alpha\beta}(A_\beta)$.}
\end{itemize}
Il est clair que si les $f_{\alpha\beta}$ sont des \emph{épimorphismes}, la
condition (ML) est vérifiée. Inversement, si (ML) est vérifiée, et si pour
tout $\alpha\in I$, $A'_\alpha$ est le sous-objet des «~images
universelles~» dans $A_\alpha$, la restriction de $f_{\alpha\beta}$ à
$A'_\beta$ est un \emph{épimorphisme}
$A'_\beta\to A'_\alpha$ pour $\alpha\leq\beta$~: en effet, si
$\gamma\geq\beta$ est tel que
$f_{\beta\delta}(A_\delta)=f_{\beta\gamma}(A_\gamma)$ pour
$\delta\geq\gamma$, on a
$A'_\beta=f_{\beta\gamma}(A_\gamma)$ et cela entraîne d'autre part
$f_{\alpha\delta}(A_\delta)=f_{\alpha\gamma}(A_\gamma)$ pour
$\delta\geq\gamma$, donc
$A'_\alpha=f_{\alpha\gamma}(A_\gamma)=f_{\alpha\beta}(A'_\beta)$.

Notons aussi que la condition (ML) est vérifiée lorsque les objets
$A_\alpha$ sont \emph{artiniens} dans $\mathcal C$, c'est-à-dire que toute
famille de sous-objets de $A_\alpha$ admet un élément minimal~: un élément
minimal de la famille filtrante décroissante
$(f_{\alpha\beta}(A_\beta))$ de sous-objets de $A_\alpha$ est en effet
alors nécessairement le plus petit de ces sous-objets.
\end{env}

\begin{remark}[13.1.3]
\label{0.13.1.3-fr}
La condition (ML) peut se formuler également lorsque $\mathcal C$ est par
exemple la catégorie des ensembles~; on peut alors encore définir le
sous-ensemble des «~images universelles~» de $A_\alpha$ et les remarques
faites à ce sujet dans (\hyperref[0.13.1.1-fr]{13.1.1}) et
(\hyperref[0.13.1.2-fr]{13.1.2}) restent valables.
\end{remark}

\subsection{La condition de Mittag-Leffler pour les groupes abéliens}
\label{subsection:0.13.2-fr}

\begin{proposition}[13.2.1]
\label{0.13.2.1-fr}
\emph{Soit}
\[
0\longrightarrow A_\alpha\xrightarrow{u_\alpha}B_\alpha
\xrightarrow{v_\alpha}C_\alpha\longrightarrow0
\]
\emph{une suite exacte de systèmes projectifs de groupes abéliens (relatifs
à un même ensemble filtrant d'indices $I$).}
\begin{enumerate}
\item[\emph{(i)}]
\emph{Si $(B_\alpha)$ vérifie (ML), il en est de même de $(C_\alpha)$.}
\item[\emph{(ii)}]
\emph{Si $(A_\alpha)$ et $(C_\alpha)$ vérifient (ML), il en est de même de
$(B_\alpha)$.}
\end{enumerate}
\end{proposition}

\begin{proof}
Soient $(f_{\alpha\beta})$, $(g_{\alpha\beta})$, $(h_{\alpha\beta})$ les
systèmes d'homomorphismes définissant les systèmes projectifs
$(A_\alpha)$, $(B_\alpha)$, $(C_\alpha)$ respectivement.

(i) Supposons que
$g_{\alpha\beta}(B_\beta)=g_{\alpha\lambda}(B_\lambda)$ pour
$\lambda\geq\beta$~; comme $v_\beta$ et $v_\lambda$ sont surjectives, on a
\[
h_{\alpha\beta}(C_\beta)
=v_\alpha(g_{\alpha\beta}(B_\beta))
=v_\alpha(g_{\alpha\lambda}(B_\lambda))
=h_{\alpha\lambda}(C_\lambda)
\]
pour $\lambda\geq\beta$.

(ii) Soit $\alpha\in I$, et soit $\beta\geq\alpha$ un indice tel que pour
$\lambda\geq\beta$, on ait
$f_{\alpha\beta}(A_\beta)=f_{\alpha\lambda}(A_\lambda)$~; soit d'autre
part $\gamma\geq\beta$ un indice tel que, pour $\lambda\geq\gamma$, on ait
$h_{\beta\gamma}(C_\gamma)=h_{\beta\lambda}(C_\lambda)$. Soit alors
$y_\alpha$ un élément de $g_{\alpha\gamma}(B_\gamma)$~; on a donc
$y_\alpha=g_{\alpha\gamma}(y_\gamma)$ avec $y_\gamma\in B_\gamma$~;
posons $y_\beta=g_{\beta\gamma}(y_\gamma)$, de sorte que
$v_\beta(y_\beta)=h_{\beta\gamma}(v_\gamma(y_\gamma))$. Pour tout
$\lambda\geq\gamma$, il existe par hypothèse $y_\lambda\in B_\lambda$ tel
que
$h_{\beta\gamma}(v_\gamma(y_\gamma))
=h_{\beta\lambda}(v_\lambda(y_\lambda))
=v_\beta(g_{\beta\lambda}(y_\lambda))$, d'où
$v_\beta(y_\beta-g_{\beta\lambda}(y_\lambda))=0$, et par suite il existe
$x_\beta\in A_\beta$
